\chapter{致谢}
在写毕业论文之前，心心念念要写致谢感慨一番，而真正写完了却提不起力气再写致谢，或是回顾这四年。有些轻舟已过万重山之感吧，无疑是我人生目前为止最重要的四年，在这里我不断成长，真正成人。

毕业论文终稿截止的那天正好是我的22岁生日，也是我写下致谢时，再向后推十小时。18岁生日时赶上成人礼，我在盘锦高中的教室里吃着蛋糕，焦虑着十天后的高考，也憧憬着大学生活。四年后的今天我坐在北京大学图书馆，在三楼东区我最熟悉的区域，窗外天阴阴的，下了点小雨，我依然焦虑未来，但我想我有了更多底气。

感谢自己，能够诚实地面对自己的感受，保持困惑地思考。感谢在最艰难的时候，也许是大一大二迷茫时，也许是大三保研时，没有放弃自己，越来越相信自己。有勇气尝试，探索更广阔的世界，再不断回到脚下前行。希望自己能不断变好，但我愿做我最忠诚的朋友。我有时感叹，自己对自己负责是最难的事情，但四年里我慢慢学会了这件事。

感谢我的父母，感谢你们永远相信我，不给我压力，鼓励和安慰一直陪伴着我，你们是最棒的父母，温暖的家永远是庇护。

感谢我的朋友们，和你们在一起永远是最开心的。生活有时候很难，看着你们也在认真过，追求想要的东西，为你们高兴，也不断鼓舞着我。喜欢与你们吃饭，闲聊，在湖边散步。

感谢一路上遇到的老师，师长，同学，你们在不同时刻给我激励，让我看到另一种可能，从你们的身上我学到很多。感谢我的导师陈语谦老师，谢谢您的认可，让我幸运地得以在深研院度过未来三年；感谢连宙辉老师，虽然和您交流不多，但第一次的科研经历让我更加坚定；感谢我的毕业论文导师俞妍老师，您的鼓励总是那么真诚，这份鼓励也会继续让我前行。在提到这些时，我总是想说，谢谢你们“收留”我，在面对并不光鲜的履历时，你们给了我机会，也许有些我最终完成得也不尽如人意，但我尝试过也努力过。谢谢一路上遇到的师长，成鹏师兄，昱博师兄，一鸣师兄，林杭师兄，宇轩师兄，模仿是最好的学习方式，我不断从你们身上汲取着能量，从零开始学会做科研，学会与人合作。感谢丰楗师兄对本毕业论文的大力支持。

感谢物理学院，虽然本科的物理课程我吸收得并不好，但学院大部分时候都让我很舒服。

感谢北大，你是最宽容的。慢慢地我发现，其实是大一第一个月我就发现，我还挺适合这里的，作为一个普通人，北大给了我无需被评判的普通的自由。喜欢园子里的草木，喜欢从近邻宝到南门高高的树，喜欢从图书馆穿过燕南园到家园食堂，喜欢图书馆三楼东区，西区301和里面的小仓库，一楼新书展阅厅，和声厅，喜欢二教207和208，三楼靠窗的自习区，喜欢理教三四楼的自习区，喜欢经院自习区，喜欢我的宿舍，33楼214和34A420，喜欢随机刷新的小猫，喜欢绕着未名湖走路，喜欢时常拥挤的理教健身房。四年过去了，我还是爱吃松林的包子，学一的襄阳牛肉面，农二的木桶饭，燕南的家常菜，家三的中式简餐，各处的麻辣香锅。习惯在闲得没事的时候去29楼和45楼地下溜达一圈，心情不好的时候去未名湖、或者从西门外走到圆明园。有时也感到压抑想要逃离，但每天晚上回来踏进校园的一刻，还是感到安心。

感谢电影，谢谢百讲，资料馆，法文，百老汇，谢谢北影节，爱酷，以及各种各样的电影活动，如果没有遇到和爱上电影我不知道会怎样。谢谢戴老师，谢谢因电影遇到的很多朋友们，能和你们相遇、一起聊天真的很幸福。大学四年“浪费时间”最多的事就是看电影了吧。喜欢/兴趣/爱好完全不需要理由，但某种程度上电影让我更加自洽。

感谢北京，你和北大一样包容，你有些破烂又无比宽广的胸襟让我见识到18岁时想要看到的世界。先是探索带来的新奇，后是厌倦，要毕业了又开始舍不得北京。不远的路我喜欢骑车，穿过街道看看两旁。我像拿着图钉标记地点，标记着那些景点，街道，小店，有些可能已经不在，但也钉在了我的回忆地图里。能在北京不同地方看到不同的建筑，不同的风格，完全不同的人做各种各样的事，因为基数和范围足够大所以有趣的事也多，我也就在探索中慢慢了解了这座城市，也许有缘再会。

感谢互联网。

继续散步吧，继续前行吧。